\section{Discussion}

The motivating result is not that one controller family is superior. It is that feedback-regulated trajectories may contain insufficient information to identify the mechanism that generated them. This has direct consequences for computational psychology: mechanistic labels can become more specific than the observations that support them.

The standardized development grid shows that weak recovery is not universal: richer families can be recovered under favorable observation conditions, while sparse or noisy observations undermine discrimination. Nominal recovery and practical predictive value are separate axes rather than interchangeable definitions of identifiability. The final numerical protocol is intended to turn these developmental boundaries into independently reproduced inferential estimates. If richer models become recoverable only under dense, low-noise sampling, that finding informs human study design. If their predictive losses remain practically equivalent, that favors simpler predictive use at the stated resolution but does not alone establish non-identifiability of the generating family. Weak family recovery and uncertain parameters require their own evidence before restricting mechanistic interpretation.

This perspective clarifies the relationship to ARC. The ARC framework does not depend on P/PI/PID or any other controller family as a psychological type. M02 asks when latent dynamical labels are recoverable at all and therefore constrains mechanistic interpretation without serving as a prerequisite for ARC.

The local Windows repeat additionally shows that archive count must not be confused with statistical replication: three fully checksum-verified collections contained the same 5,760 trajectory records. The targeted aggregate nominal recovery (51.9\%) and trajectory-level near-tie fraction (90.4\% for the developmental tolerance) are descriptive, not multiplicity-controlled evidence for equivalence. They highlight the need to differentiate recovery, small predictive contrasts, and formal interval-contained practical equivalence.

Limitations include dependence on the chosen generator set, practical-equivalence tolerance, simulator dynamics, and observation process. A Monte Carlo study can expose ambiguity within the investigated model universe but cannot prove that all plausible mechanisms have been considered.


\subsection{Relational Motivation and Measurement Implications}

The application to relationships motivates observation design rather than validating a controller interpretation. In a dyad, both partners act and influence each other's subsequent state; the present single-plant feedback coefficient aggregates environmental response and does not estimate reciprocal co-regulation. The following implications are hypotheses for measurement design, not observed human findings.

\begin{table}[htbp]
\centering
\caption{Relational hypotheses motivating future observation design}
\label{tab:m02-relational-implications}
\begin{tabular}{p{0.21\linewidth}p{0.33\linewidth}p{0.33\linewidth}}
\toprule
Candidate process & Observation consequence & Measurement implication \\
\midrule
Reciprocal recovery & Stable feedback may narrow the range of observable discrepancy & Observe both partners and external disturbances; low variation need not mean weak regulation \\
Brief reactivity & Responses may occur between widely spaced prompts & Record timing and event-linked observations; sparse samples can alias fast changes \\
Unobserved regulatory channels & Similar recorded communication can coexist with different internal or overt actions & Separate perceived state, desired state and action channels; low observed action does not identify emotional suppression \\
\bottomrule
\end{tabular}
\end{table}

These possibilities explain why richer ecological measurements may be necessary, but a favorable synthetic design cannot guarantee psychological identifiability. M03 must establish defensible measurement, and M04 must test prospective behavioral associations independently. Neither manuscript should treat fitted PID gains as validated attachment traits. Full dyadic likelihoods and genuinely ordinal measurement models are separate extensions, not properties of the current benchmark.
